%% =========================================================================
\section{Introduction}
%% =========================================================================

This report covers the simulation of a three-phase synchronous generator using
MATLAB R2025b and Simscape Electrical. There are three exercises with six
tasks, each one building on the previous.

Exercise~1 looks at how a \MVA{555}, \kV{24}, \Hz{60} round-rotor generator
responds to step load changes. It covers the relationship between active power
and system frequency (through the swing equation) and between reactive power
and terminal voltage (through the AVR). Exercise~2 extends this to two
identical generators running in parallel. The governor frequency setpoint
controls how load is shared between them --- an under-frequency machine absorbs
power and acts as a motor, while an over-frequency machine picks up the load.
The system frequency settles at a droop-weighted average of the two setpoints.
Exercise~3 swaps one generator for an ideal infinite bus. The infinite bus
holds frequency fixed at \pu{1.0}, so the oncoming generator has no effect on
system frequency. Its power output depends only on its governor setpoint
relative to the grid frequency. Simulation plots, equation derivations, and
steady-state calculations support the results throughout.
